\documentclass[aspectratio=169]{beamer}

% -------------------------------------------------
% Pacotes
% -------------------------------------------------
\usepackage[utf8]{inputenc}
\usepackage[brazil]{babel}
\usepackage{graphicx}
\usepackage{xcolor}
\usepackage{hyperref}
\usepackage{float}

% -------------------------------------------------
% Tema e cores
% -------------------------------------------------
\usetheme{Madrid}
\usecolortheme{seahorse}
\usefonttheme{professionalfonts}

% -------------------------------------------------
% Título
% -------------------------------------------------
\title[Tomada de Decisão Baseada na Evidência]{Apontamento de Literatura – Tema 2\\\large Introdução à Tomada de Decisão Baseada na Evidência e às Revisões Sistemáticas}
\author[Disciplina]{Disciplina: Revisão Sistemática e Meta-análises\\Pós-graduação (Mestrado e Doutoramento)}
\date{}

% -------------------------------------------------
\begin{document}

\maketitle

% =================================================
\begin{frame}{Evidência: conceito fundamental}
O termo \textbf{evidência} refere-se a informação empírica ou teórica utilizada para sustentar uma decisão. É \textbf{conhecimento provisório}, sujeito a revisão, crítica e falseamento.

\textbf{Características da evidência científica:}
\begin{itemize}
    \item Obtida por métodos explícitos e sistemáticos
    \item Reprodutível ou verificável
    \item Sujeita à revisão por pares
    \item Possui incerteza quantificável
\end{itemize}

\begin{figure}
    \centering
    \includegraphics[width=0.6\textwidth]{fig1_evidencia.png}
    \caption{Diagrama conceptual da evidência científica}
\end{figure}
\end{frame}

% =================================================
\begin{frame}{Evidência científica}
A evidência científica resulta da aplicação do método científico a fenómenos naturais ou sociais. Distingue-se de:
\begin{itemize}
    \item Evidência anedótica
    \item Opiniões de especialistas não sistematizadas
    \item Experiência pessoal não controlada
\end{itemize}

Em termos epistemológicos, a evidência científica é:
\begin{itemize}
    \item \textbf{Empírica} (dados observáveis)
    \item \textbf{Teórica} (interpretada à luz de modelos e hipóteses)
    \item \textbf{Probabilística} (nunca absolutamente certa)
\end{itemize}
\end{frame}

% =================================================
\begin{frame}{Práticas Baseadas na Evidência (EBP)}
Integra sistematicamente:
\begin{enumerate}
    \item Melhor evidência científica disponível
    \item Experiência profissional especializada
    \item Contexto, valores e necessidades dos utilizadores finais
\end{enumerate}

\begin{figure}
    \centering
    \includegraphics[width=0.5\textwidth]{fig2_ebp.png}
    \caption{Triângulo da prática baseada na evidência (EBP)}
\end{figure}
\end{frame}

% =================================================
\begin{frame}{Necessidade de falseabilidade}
Segundo Karl Popper, uma afirmação científica deve ser \textbf{falseável}.

\begin{figure}
    \centering
    \includegraphics[width=0.7\textwidth]{fig3_falseabilidade.png}
    \caption{Fluxo de hipótese, teste e falseamento}
\end{figure}
\end{frame}

% =================================================
\begin{frame}{Fluxo da evidência: da geração à implementação}
Etapas do ciclo da evidência:
\begin{enumerate}
    \item Geração da evidência
    \item Síntese da evidência
    \item Transferência do conhecimento
    \item Implementação prática
    \item Avaliação e retroalimentação
\end{enumerate}

\begin{figure}
    \centering
    \includegraphics[width=0.8\textwidth]{fig4_fluxo_evidencia.png}
    \caption{Fluxo da evidência desde a geração até à implementação}
\end{figure}
\end{frame}

% =================================================
\begin{frame}{O ecossistema da evidência}
Atores principais: investigadores, revisores, instituições, sociedade civil e decisores.

\begin{figure}
    \centering
    \includegraphics[width=0.6\textwidth]{fig5_ecossistema.png}
    \caption{Rede do ecossistema da evidência}
\end{figure}
\end{frame}

% =================================================
\begin{frame}{Síntese da evidência}
A síntese da evidência é necessária para reduzir incertezas e evitar decisões baseadas em estudos isolados.

\begin{figure}
    \centering
    \includegraphics[width=0.5\textwidth]{fig6_sintese.png}
    \caption{Pirâmide da síntese da evidência}
\end{figure}
\end{frame}

% =================================================
\begin{frame}{Tipos de revisão}
\begin{itemize}
    \item \textbf{Revisão narrativa}: não sistemática, subjetiva, contextualização teórica
    \item \textbf{Revisão de escopo}: mapeamento de área e lacunas
    \item \textbf{Revisão rápida}: metodologia simplificada, decisões urgentes
    \item \textbf{Revisão sistemática}: metodologia explícita, avaliação crítica, síntese qualitativa e quantitativa
\end{itemize}

\begin{figure}
    \centering
    \includegraphics[width=0.6\textwidth]{fig7_tipos_revisao.png}
    \caption{Comparação entre tipos de revisão}
\end{figure}
\end{frame}

% =================================================
\begin{frame}{Revisão sistemática}
Principais elementos:
\begin{itemize}
    \item Questão estruturada (PICO)
    \item Protocolo pré-definido
    \item Critérios de inclusão/exclusão
    \item Busca abrangente
    \item Avaliação da qualidade
    \item Síntese transparente
\end{itemize}

\begin{figure}
    \centering
    \includegraphics[width=0.8\textwidth]{fig8_fluxo_revisao.png}
    \caption{Fluxograma do processo de revisão sistemática}
\end{figure}
\end{frame}

% =================================================
\begin{frame}{Preparação para uma revisão sistemática}
Checklist:
\begin{itemize}
    \item Verificar revisões existentes
    \item Avaliar viabilidade
    \item Definir objetivos e escopo
    \item Escolher diretrizes metodológicas
    \item Registo no PROSPERO
\end{itemize}

\begin{figure}
    \centering
    \includegraphics[width=0.6\textwidth]{fig9_checklist.png}
    \caption{Checklist de preparação para revisão sistemática}
\end{figure}
\end{frame}

% =================================================
\begin{frame}{Processo completo de revisão sistemática}
O processo completo envolve: questão → protocolo → registo → busca → seleção → extração → avaliação → síntese → reporte.

\begin{figure}
    \centering
    \includegraphics[width=0.9\textwidth]{fig10_processo_revisao.png}
    \caption{Fluxo completo da revisão sistemática}
\end{figure}
\end{frame}

% =================================================
\begin{frame}{Considerações finais}
Para mestrados e doutoramentos, a revisão sistemática é:
\begin{itemize}
    \item Uma postura epistemológica
    \item Um instrumento de decisão científica
    \item Um pilar da ciência moderna baseada na evidência
\end{itemize}
\end{frame}

\end{document}
